\chapter{開発するアプリケーション }

\section{目的とコンセプト}
公立はこだて未来大学での生活は不便なことが多々ある．大学が提供しているHOPEやSTUDENTSは情報が散乱していてどこにどの情報があるのかすぐにわからないという問題がある．そのため私達は，公立はこだて未来大学に在籍している学生に向けてHOPEやSTUDENTSにある情報をまとめるアプリを開発し，公立はこだて未来大学の学生が簡単に欲しい情報にアクセスできるようにすることを目的とした．そして，システムに存在する情報だけでなく，公立はこだて未来大学に在籍している学生として生活に必要な情報もまとめることで大学生活の利便性を高めることが目標である．

\bunseki{大須賀雅也}

\section{機能}
私達は，公立はこだて未来大学での生活において不便である点を考え，それらを解消できるような機能についてのアイディア出しを行った．既存の公立はこだて未来大学のシステムで使いにくい部分を解消する機能に加えて，「こんな機能があったらいいな」という機能を中心にアイディアを出した．それらに基づいて，実装を検討している機能を紹介する．
\bunseki{大須賀雅也}

\subsection{授業系の情報をまとめる機能}
この機能では，シラバスや学生の口コミなどの科目情報をまとめ，見やすく表示する．シラバスについては中間試験・期末試験の有無やレポートの有無などの成績基準の情報や，必修であるか選択であるかなどの区分を詳しく表示する．また，科目ごとに学生の口コミを投稿できるフォームを用意する．これらを表示することで，学生の履修登録時の科目選択の参考にしてもらうことを目標としている．
\bunseki{南川虎之介}

\subsection{バイトやイベントなどの情報をまとめる機能}
この機能では，学生または教員がバイト募集やイベントなどの告知をできる．これらの情報は学内メールや大学構内の掲示板，学生個人のSNSなどに分散している．これらの情報をこの機能内で投稿してもらい，ユーザがこの機能だけで情報を収集できるようにする．
\bunseki{南川虎之介}

\subsection{空き教室やテーブルの情報を表示する学内マップ機能}
この機能では，空き教室やテーブルの情報を表示する学内マップを表示することができる．STUDENTSから教室の予約情報を取得し，空き教室を表示する．テーブルの情報とは，コンセントの位置や席数など，学生が席を探すうえで必要な情報を表示する．この機能により，学生が教室や席を探す際の負担を軽減し，移動にかかる時間を減らすことができる．
\bunseki{大須賀雅也}

\subsection{学生同士で不必要な物品をやり取りする機能}
この機能では，不要になったがまだ利用価値があるものを受け渡すことができる．例えば，教科書・車・家電等の受け渡しの仲介をすることで, 必要なものをお金が無くて買えない学生の経済的な負担を軽減する. 不要なものを写真や詳細付きで投稿でき, それを欲しいユーザに条件が合えば受け渡すことができる. 
\bunseki{齊藤輝}

\subsection{過去問をまとめる機能}
この機能ではユーザが中間試験や期末試験等の過去問のダウンロードやアップロードを自由にできる．コースや学年などを選択し，それにより絞り込まれた科目一覧が表示される．科目名をタップすると年度ごとの過去問が表示される．先輩との関わりがない学生をメインターゲットとしている．
\bunseki{齊藤輝}

